Em políticas públicas e avaliações de programas, um dos desafios centrais é estimar efetivamente qual o impacto de uma intervenção -- ou seja, distinguir o que mudou \textbf{por causa} da política do que teria mudado de qualquer forma. Para fazer isso com rigor, precisamos lidar com os chamados \textit{back-doors} no modelo causal: variáveis que influenciam tanto a participação na intervenção quanto o resultado e, assim, ameaçam a validade da comparação. 

Num experimento, onde os grupos de tratamento e controle são aleatorizados, essas variáveis não seriam um problema para nossa comparação, já que a aleatorização da alocação do tratamento nos garantiria que as características dos dois grupos seriam iguais em média. No entanto, nem sempre conseguimos realizar sorteios, e selecionar todos os não-beneficiários de um programa como grupo de comparação provavelmente não resultará em uma comparação crível. 

Considere o exemplo de um programa de qualificação profissional para jovens de baixa renda, como cursos de curta duração voltados para inserção no mercado de trabalho. Se compararmos diretamente a taxa de ocupação de participantes e não participantes, talvez encontremos que os primeiros têm maior empregabilidade após o curso.  Mas esse resultado pode ser enganoso: é possível que os jovens que se inscreveram já fossem, em média, mais motivados, tivessem redes de contato melhores, ou nível educacional mais alto que os não inscritos. Nesse caso, a diferença observada não refletiria o efeito do programa, mas sim essas diferenças pré-existentes — o chamado \textbf{viés de seleção}.  

Esse problema é ubíquo em políticas públicas. Em saúde, por exemplo, indivíduos que aderem a campanhas de prevenção podem já ser mais conscientes sobre hábitos saudáveis; em educação, famílias que matriculam filhos em escolas de tempo integral podem diferir sistematicamente de outras em termos de renda ou engajamento. 

Por isso, o desafio central da avaliação é criar comparações justas, isto é, contrafactuais\footnote{Contrafactuais representam os resultados potenciais que uma unidade teria apresentado caso estivesse submetida a uma condição diferente da observada. O efeito causal é definido como a diferença entre o resultado observado e o resultado contrafactual, este último não observável e estimado por métodos estatísticos \citep{ImbensRubin2015}.} que representem de forma plausível o que teria acontecido com os beneficiários caso não tivessem recebido o tratamento. Em termos operacionais, comparações justas exigem que unidades tratadas e não tratadas sejam comparáveis em todas as covariáveis relevantes que afetam simultaneamente a probabilidade de tratamento e o resultado, de modo que, condicional a essas características observáveis, a diferença média de resultados possa ser interpretada como efeito causal do programa.

O pareamento oferece uma resposta intuitiva a esse desafio: \textbf{para avaliar um programa, precisamos encontrar um grupo de comparação que se pareça o máximo possível com o grupo tratado}. Importante notar, contudo, que o pareamento não cria contrafactuais perfeitos, mas sim aproximações empíricas que são válidas sob hipóteses específicas sobre o processo de seleção. Em particular, o método procura simular uma situação experimental na qual tratados e controles são semelhantes em todas as características observáveis relevantes, diferenciando-se apenas pelo fato de terem (ou não) participado da política.


Voltando ao exemplo dos cursos de qualificação: em vez de comparar todos os inscritos com todos os não inscritos, podemos restringir a comparação apenas a jovens com perfis muito semelhantes — por exemplo, mesma faixa etária, nível de escolaridade, gênero e bairro de residência. Essa estratégia, no entanto, só é viável quando existe \textbf{suporte comum} - quando para cada perfil observado entre os participantes do curso há indivíduos comparáveis no grupo de não inscritos. Sob essa condição, se o curso teve impacto real, a diferença nos resultados de emprego entre esses grupos comparáveis será uma estimativa mais próxima do efeito verdadeiro.


A intuição é simples: ao eliminar diferenças observáveis entre tratados e controles, reduzimos o risco de atribuir ao programa algo que se deve, na verdade, a características pré-existentes dos participantes. É por isso que o \textit{pareamento} é frequentemente descrito como um método de \textbf{``construir contrafactuais''} usando dados observacionais. Essa estratégia, contudo, é suficiente apenas na medida em que as principais fontes de seleção no programa sejam observáveis e adequadamente controladas, hipótese conhecida como \textit{ignorabilidade} ou seleção em observáveis.


Essa abordagem é especialmente útil quando há um número razoável de variáveis observáveis que influenciam tanto a participação no tratamento quanto o desfecho de interesse, e quando há sobreposição entre os grupos. Assim, o pareamento se apresenta como uma ferramenta poderosa para estudos aplicados que buscam inferência causal em ambientes não experimentais.

Este guia tem como objetivo apresentar de forma clara e sistemática os principais conceitos, métodos e aplicações do pareamento em econometria, destacando tanto a intuição quanto os aspectos técnicos das abordagens.

O pareamento, apesar de ser uma ferramenta poderosa, requer cuidados na presença de viés de seleção associado a características observáveis, bem como na escolha de métricas de distância e na avaliação da qualidade do balanceamento entre tratados e controles. Esses aspectos dizem respeito à capacidade do método de eliminar diferenças sistemáticas em covariáveis observadas. Por outro lado, vieses decorrentes de fatores não observáveis que afetam simultaneamente a participação no programa e os resultados não são, em geral, corrigíveis por pareamento. Portanto, também discutiremos técnicas de diagnóstico e procedimentos de avaliação da qualidade do pareamento, conforme sugerido por \cite{abadie2006large} e \cite{imbens2004nonparametric}.

O restante do guia está organizado da seguinte forma. A Seção 2 apresenta o arcabouço causal, introduzindo o modelo de resultados potenciais, o problema do viés de seleção e a distinção entre dados experimentais e não-experimentais, utilizando um programa real como exemplo motivador. A Seção 3 introduz o método de pareamento, discutindo sua lógica de identificação e sua relação com a construção de contrafactuais em dados observacionais. As Seções 4 e 5 detalham os principais métodos operacionais de pareamento, começando pelo caso univariado e avançando para o pareamento com múltiplas covariáveis, incluindo o uso do \textit{propensity score}. Por fim, a Seção 6 apresenta as conclusões e discute limitações e extensões do método.